\section{Introduction}

The Auslander--Reiten conjecture asks whether a finitely generated module
$M$ over an Artin algebra $R$ must be projective when
\begin{equation}\label{eq:arc}
 \Ext_R^i(M,M\oplus R)=0\qquad\text{for every }i>0.
\end{equation}
Here an Artin algebra is finite over a commutative Artinian central base
ring. In particular, every finite-dimensional algebra over a field is
an Artin algebra. The conjecture originates in the work of Maurice Auslander and
Idun Reiten on the generalized Nakayama conjecture
\cite{AuslanderReiten1975}; its module and generator formulations are
recalled in \cite[Introduction]{HunekeLeuschke2004}.

We give a counterexample. A module is called \emph{Gorenstein-projective}
if it is a cokernel in a doubly infinite exact complex of finite
projectives that remains exact after applying $\Hom_R(-,R)$.
The \emph{Gorenstein-projective conjecture} asks whether a finite
Gorenstein-projective module $M$ with $\Ext_R^i(M,M)=0$ for every
$i>0$ must be projective~\cite[p.~2]{LuoHuang2008}. The ambient algebra
is not assumed to be Gorenstein. Total acyclicity gives
$\Ext_R^i(M,R)=0$ for $i>0$, so this conjecture is precisely the
restriction of the Auslander--Reiten conjecture to Gorenstein-projective
modules.

\begin{theorem}\label{thm:main}
Let $\kk=\mathbb F_2(q,H_1,H_2)$, where $q,H_1,H_2$ are algebraically
independent. There exist a finite-dimensional $\kk$-algebra $\Lambda$
and a finite-dimensional nonprojective left $\Lambda$-module $Z$ such that
\[
 \Ext^i_\Lambda(Z,Z)=0=\Ext^i_\Lambda(Z,\Lambda)
 \qquad\text{for every integer }i\geq1.
\]
The module $Z$ is Gorenstein-projective. The semisimple quotient is
$\Lambda/\rad\Lambda\simeq\kk^8$, so the algebra is split with eight
isomorphism classes of simple modules; also $\rad(\Lambda)^4\ne0$.
These conclusions persist after every field extension of $\kk$.
\end{theorem}

Thus both conjectures fail for Artin algebras, including finite-dimensional
algebras over an algebraically closed field of characteristic two:
one can extend scalars to an algebraic closure of $\kk$. The parameters
play essential roles in the proof. The powers of $q$ give an infinite
resolution with uniform kernel formulas, and $H_1/H_2$ has infinite
multiplicative order. The algebraic closure in this assertion contains
these transcendental elements.

The final algebra is triangular and noncommutative. The companion article
\cite[Theorem~1.1 and Corollary~1.2]{OpenAITachikawa2026} constructs a
symmetric counterexample over the
same field by a different bimodule construction and a transfer requiring
control of negative complete self-extensions, and develops the
homological consequences of its associated endomorphism algebra.
The contribution here is
the explicit eight-simple Gorenstein-projective example and its persistence
under every field extension. Each paper contains its own complete proof.
The commutative restriction is not resolved by either construction; see
\cite[Appendix A]{ChristensenHolm2010} for its relationship with the
various Auslander and Tachikawa conditions.
The generator formulation also fails: $Z\oplus\Lambda$ is a
nonprojective generator with no positive self-extensions. Indeed Ext is
additive in both arguments, and all positive Ext groups with first
argument $\Lambda$ vanish.

\subsection{Earlier constructions and restricted results}

The positive cases developed along several lines. For the commutative
Noetherian formulation, Auslander, Ding and Solberg treated complete
local complete intersections by lifting across
regular sequences \cite[Proposition~1.9]{AuslanderDingSolberg1993}.
Huneke and Leuschke proved the conjecture for excellent commutative
Cohen--Macaulay normal domains containing the rational numbers, among
other classes \cite[Theorem~0.1]{HunekeLeuschke2004}.
Christensen and Holm established a general cohomological criterion:
every left-noetherian ring satisfying Auslander's condition satisfies
the conjecture \cite[Theorem~A]{ChristensenHolm2010}. That condition
requires, for each fixed finitely generated first module, a common bound
after which Ext vanishes for every finitely generated second module whose Ext groups
vanish eventually.

Algebras with radical cube zero form another important positive class.
Huneke, \c{S}ega and Vraciu treated commutative Artinian local rings
\cite[Theorem~4.1(1)]{HunekeSegaVraciu2004}. For split finite-dimensional
local algebras with radical cube zero, Ringel and Zhang proved the
stronger assertion that every nonprojective module with no positive
Ext to the algebra has a nonzero first self-extension
\cite[Theorem~1.5]{RingelZhang2022}.
Zhang and Zhou subsequently proved the conjecture for split
finite-dimensional algebras with radical cube zero
\cite[Theorem~1.1]{ZhangZhou2026}. Their result for Igusa--Todorov
Artin algebras includes all Artin algebras with radical cube zero
\cite[Theorem~3.2 and Corollary~4.4]{ZhangZhouIgusaTodorov2026}.
The algebra constructed here has nonzero fourth radical power, as
proved in Proposition~\ref{prop:radical}.

Schulz's construction of a nonprojective module without self-extensions
\cite{Schulz1994} is an important antecedent, but its particular
skew-field example lies outside the Artin-algebra setting above; see
\cite[Introduction]{DiverisPurin2012}. A separate distinction is between
all-degree and eventual vanishing: both summands in \eqref{eq:arc}
must vanish in every positive degree.

Schulz also exhibited two-dimensional modules over quantum exterior
algebras whose syzygies move through a parameter family and whose
self-extensions vanish eventually, while the first self-extension
remains nonzero \cite[Example~7]{Schulz1995}. The corresponding quantum
exterior algebra occurs as a corner of our starting algebra. We use
the same parameter-shifting mechanism for a resolution over a larger
algebra, then remove every positive self-extension by two independent twists.
For indecomposable weakly symmetric algebras over an algebraically closed
field with $J^3=0\ne J^2$, Erdmann characterizes those admitting
nonprojective modules with eventually vanishing self-extensions
\cite{Erdmann2017}; eventual vanishing alone does not give
\eqref{eq:arc}.
Xia proves the Auslander--Reiten conjecture for quantum complete
intersections \cite{Xia2026}, including the quantum exterior algebras
used in Schulz's parameter-shifting examples.

Vanishing in a finite range also differs from the hypothesis of
\eqref{eq:arc}. B\"ohmler and
Marczinzik \cite[Theorem 3.1 and the end of Section 3]{BoehmlerMarczinzik2026}
construct a characteristic-two commutative self-injective example with
vanishing first and second self-extensions, and nonzero third
self-extension.

Tachikawa's second conjecture is the self-injective case of
\eqref{eq:arc}. Enomoto proves that this conjecture for
the two-fold trivial extension of an Artin algebra implies the
Auslander--Reiten conjecture for that algebra
\cite[Theorem~3.1]{Enomoto2026}. His reduction produces a self-injective
algebra. The symmetric construction
in our companion and the structural assertions of Theorem~\ref{thm:main}
are separate contributions.

\subsection{How the construction removes self-extensions}

The proof has two parts. First we give a conversion principle: a
suitable stable map $v:X\to F\otimes_A X$ over a symmetric algebra
$A$ produces a module over
\[
 \Lambda=\begin{pmatrix}A&0\\F&A\end{pmatrix}.
\]
Its positive self-extensions are controlled by the difference between
the actions of a self-extension of $X$ on the two sides of $v$.
Total acyclicity separately gives the vanishing of Ext into $\Lambda$.
This principle is stated and proved before the example, so the
subsequent construction has an explicit target.

We then build the required $A,F,X,v$. Write $D$ for the $\kk$-linear dual.
A ten-dimensional algebra $C$
has a simple module $s$ whose Ext groups are especially sparse. Passing
to its symmetric trivial extension $T=C\ltimes DC$ produces
$\Ext_T^*(s,s)=\kk[\tau]$ with $|\tau|=3$. An explicit Hochschild
cocycle realizes $\tau$ and supplies a homotopy compatible with twists
of $T$. For $A=T\otimes T$ and $X=s\otimes s$, the nonnegative
self-extension algebra is $\kk[\tau_1,\tau_2]$.

Write $\mathcal C$ for the bimodule obtained from the two cones of
these classes. In the stable category, maps factoring through projective
modules are identified with zero, and the shift $[1]$ is inverse syzygy.
The stable maps from $X$ to
$(\mathcal C\otimes_A X)[a]$ form a line for $a=0$ and $a=3$ and
vanish in every other integer degree. This profile supplies a common
target for two differently twisted maps.

For a nonzero scalar $H$, let $h_H$ fix $C$ and scale $DC$ by $H$,
and let $U_H$ be $A$ with right action twisted by $h_H\otimes h_H$.
Its tensor functor fixes $X$ but multiplies its degree-$3m$
self-extensions by $H^{-m}$. Using the symmetrizing trace and the
cocycle's boundary identity, we construct maps $U_H\to\mathcal C[3]$
whose evaluations on $X$ are nonzero. After rescaling the maps for
$H_1$ and $H_2$, their evaluations are the same nonzero element of
that degree-three line.
Their finite fiber $F$ then has a single diagonal stable map
$v:X\to F\otimes_A X$ in degree zero, with both projections equal to
$\id_X$. In positive degrees, projection identifies the stable maps
from $X$ to $(F\otimes_A X)[a]$ with two copies of the self-extensions
of $X$. In degree $3m>0$, the comparison takes two independently chosen
self-extension classes $g,j$ to $F(g)v+v[3m]j$, where $F(g)$ denotes
tensoring with $F$. Under these projections its matrix is
\[
 \begin{pmatrix}H_1^{-m}&1\\H_2^{-m}&1\end{pmatrix}
 \otimes\id_{\kk[\tau_1,\tau_2]_{3m}}.
\]
This matrix is invertible for every $m>0$. In degree zero, the
comparison is surjective with a nonzero kernel. The conversion
principle therefore removes all positive self-extensions while keeping
a nonzero stable endomorphism, which proves nonprojectivity.

\begin{figure}[ht]
\centering
\begin{tikzpicture}[>=Stealth,every node/.style={font=\small},
 box/.style={draw,rounded corners=2pt,inner sep=6pt,align=center},
 lab/.style={font=\footnotesize,align=center}]
\node[box] (c) at (0,3.3) {$C$\\$\dim C=10$};
\node[box] (t) at (4.2,3.3) {$T=C\ltimes DC$\\$\Ext_T^*(s,s)=\kk[\tau]$};
\node[box] (a) at (8.6,3.3) {$A=T\otimes T$\\$X=s\otimes s$};
\node[box] (f) at (0,1.4) {$F$};
\node[box] (u) at (4.2,1.4) {$U_{H_1}\oplus U_{H_2}$};
\node[box] (cone) at (8.6,1.4) {$\mathcal C[3]$};
\node[box] (lam) at (0,-.3) {$\Lambda=\left(\begin{smallmatrix}A&0\\F&A\end{smallmatrix}\right)$};
\node[box] (z) at (4.2,-.3) {$Z$ from $X,v$};
\draw[->] (c)--node[above=9mm,lab]{trivial extension}(t);
\draw[->] (t)--node[above=9mm,lab]{tensor square}(a);
\draw[->] (a)--node[right,lab]{two cones\\and shift}(cone);
\draw[->] (f)--node[above,lab]{fiber projection}(u);
\draw[->] (u)--node[above,lab]{lifted twists}(cone);
\draw[->] (f)--node[left,lab]{triangular\\algebra}(lam);
\draw[->] (lam)--node[above,lab]{finite quotient}(z);
\end{tikzpicture}
\caption{The algebraic and homological stages of the construction.
The middle row begins the stable fiber triangle; the other arrows indicate
construction steps. The finite bimodule $F$ is the ordinary kernel in
\eqref{branch:finite-fiber} before $\Lambda$ is formed. Evaluating this
fiber on $X$ gives the diagonal stable map $v$; the final module is
the finite quotient \eqref{conv:finite-module}, using both $X$ and $v$.}\label{fig:construction}
\end{figure}

The main construction-specific ingredients are the explicit cocycle
with its twisted boundary identity, the lift through both cones, and
the resulting uniform comparison of the two twist actions. The
surrounding tools have a substantial history. Trivial extensions and
their Hochschild complexes are studied in
\cite{CibilsMarcosRedondoSolotar2003}; Hochschild evaluation as a
characteristic action is treated in \cite[Section~4.3]{Lowen2008}.
Bimodules attached to Hochschild classes and their iterated tensor products
occur in the support-variety constructions of Erdmann, Holloway,
Snashall, Solberg and Taillefer \cite[Section~2]{ErdmannEtAl2003}.
We use the natural symmetric stable duality of
\cite[Section 2]{Linckelmann2013} and the complete-resolution methods
of \cite{Veliche2006}. Gorenstein-projective modules over triangular
matrix rings are studied in \cite{EshraghiEtAl2014}. We supply the
particular comparison and finiteness arguments needed here.

For reference, the objects and their roles are summarized in
Table~\ref{tab:objects}. The finite cocycle table appears in
Appendix~\ref{app:cochain}; its identities are finite algebra
calculations. The resolution, duality, cone, and complete-resolution
arguments establish the assertions in all degrees.

Section~\ref{sec:stable} fixes the stable-category conventions, and
Section~\ref{conv:section} proves the conversion principle.
Sections~\ref{alg:section} and~\ref{res:section} give the starting
algebra and its self-extension calculation.
Section~\ref{cone:section} computes the two-cone profile;
Sections~\ref{sec:lift} and~\ref{branch:section} construct the lifts and
the finite fiber. Section~\ref{sec:consequences} proves the structural
and scalar-extension assertions.

\begin{table}[ht]
\centering\small
\begin{tabular}{@{}lll@{}}
\toprule
Object & Ambient algebra & Role or dimension\\
\midrule
$C$ & $\kk$ & ten-dimensional algebra\\
$T=C\ltimes DC$, $s$ & $\kk$, $T$ & symmetric algebra of dimension $20$; simple\\
$A=T\otimes T$, $X=s\otimes s$ & $\kk$, $A$ & symmetric algebra of dimension $400$; simple\\
$\mathcal C$ & stable bimodules over $A$ & two cones of degree-three classes\\
$U_{H_i}$ & $A$-bimodules & right twists of $A$\\
$F$ & $A$-bimodules & finite fiber, projective on both sides\\
$v:X\to F\otimes_A X$ & stable $A$-modules & both projections equal $\id_X$\\
$\Lambda$, $Z$ & $\kk$, $\Lambda$ & dimension $800+\dim_\kk F$; finite module\\
\bottomrule
\end{tabular}
\caption{Objects used in the proof. Shifts are homological.}
\label{tab:objects}
\end{table}
